\documentclass[11pt]{article}

\usepackage[margin=1in]{geometry}
\usepackage{graphicx}
\usepackage{float}
\usepackage{amsmath}
\usepackage{amssymb}
\usepackage{booktabs}
\usepackage{xcolor}
\usepackage{hyperref}
\usepackage{enumitem}
\usepackage{caption}
\usepackage{tikz}
\usetikzlibrary{arrows.meta, positioning, fit, backgrounds, calc, shapes.geometric}

\hypersetup{colorlinks=true, linkcolor=blue, urlcolor=blue, citecolor=blue}

\title{\textbf{Proprioception-Cued Object Identification\\ for Learning from Demonstration}\\[4pt]
\large Progress writeup --- carrot pick/press/drop trial (\texttt{lfdws\_t001})}
\author{Anurag Shetye \quad\textbar\quad UTwente / NAKAMA Robotics Lab (remote)}
\date{\today}

\begin{document}
\maketitle

\section{Goal}

Given a demonstration recorded as a ROS\,2 bag (robot pose, gripper state,
force/torque, and ZED RGB camera), automatically determine \emph{which objects
in the scene are task-relevant} --- without any manual labelling. The central
idea: the robot's own proprioception tells us \emph{when} something matters
(a grasp, a contact, a release), and those moments tell the vision system
\emph{where} to look in the corresponding camera frame.

This document records the current setup, the data and code used, and a first
result on the carrot pick/press/drop trial, followed by a prospective research
framing.

\section{Setup and Data}

The lab exports each bag with \texttt{ros2\_unbag} plus an in-house merge
script, producing a single synchronised CSV resampled onto the robot-pose clock,
alongside a folder of extracted camera frames. We build directly on that output;
no custom export code is maintained on our side.

\paragraph{Trial.} \texttt{lfdws\_t001} --- the robot picks a carrot from a
plate (gripper close), presses it onto a cup (force contact), and drops it
(gripper open). Duration $\approx 33$\,s, four topics:

\begin{table}[H]
\centering
\begin{tabular}{llr}
\toprule
\textbf{Topic} & \textbf{Type} & \textbf{Messages} \\
\midrule
\texttt{/NS\_1/.../current\_pose}        & PoseStamped       & 33{,}137 \\
\texttt{/bota\_post/wrench\_body\_compensated} & WrenchStamped & 17{,}047 \\
\texttt{/NS\_1/franka\_gripper/joint\_states}  & JointState    & 472 \\
\texttt{/zed/.../image/compressed}       & CompressedImage   & 497 \\
\bottomrule
\end{tabular}
\caption{Topics in the \texttt{lfdws\_t001} bag.}
\end{table}

\paragraph{Key input files.}
\begin{itemize}[noitemsep]
  \item \texttt{lfdws\_t001\_0.csv} --- merged, pose-synchronised table
        (33{,}137 rows). Columns are dot-separated, topic-prefixed; the final
        column references the ZED frame by timestamp.
  \item \texttt{zed\_zed\_node\_rgb\_color\_\allowbreak rect\_image\_compressed/}
        --- 497 PNG frames, each named
        \texttt{\{nanosecond\_\allowbreak timestamp\}.png}. The integer in the
        CSV's image column is exactly the filename.
\end{itemize}

\section{Code}

A single analysis script reads the merged CSV, detects interaction events from
the proprioceptive signals, and pulls the camera frame at each event.

\paragraph{\texttt{analyze\_demo.py}.} Stages:
\begin{enumerate}[noitemsep]
  \item \textbf{Load \& parse.} Read the merged CSV; parse the gripper
        \texttt{position} column (a string such as \texttt{"[0.031, 0.031]"})
        into a scalar gripper width; compute force magnitude
        $|F|=\sqrt{f_x^2+f_y^2+f_z^2}$; build a relative-time axis from the
        pose timestamps.
  \item \textbf{Event detection.} The gripper width is near-bimodal
        (open $\approx 0.063$\,m, closed $\approx 0.003$\,m); events are found by
        threshold crossings at the midpoint (open$\to$closed = grasp,
        closed$\to$open = release). The force-contact event is the peak of the
        baseline-subtracted force magnitude restricted to the held-object window.
  \item \textbf{Timeline plot.} Gripper width and force magnitude on a shared
        time axis with event markers (Figure~\ref{fig:timeline}).
  \item \textbf{Event frames.} For each event, look up the image timestamp in
        that row and load the matching PNG (Figure~\ref{fig:frames}).
\end{enumerate}

\paragraph{Environment.} Pure offline Python (\texttt{pandas}, \texttt{numpy},
\texttt{matplotlib}, \texttt{Pillow}); no ROS runtime needed since we work from
the exported CSV and PNGs.

\section{Result}

The pick/press/drop structure is cleanly recoverable from proprioception alone,
and the camera frames at the detected event timestamps contain the correct
task-relevant objects.

\begin{table}[H]
\centering
\begin{tabular}{llll}
\toprule
\textbf{Event} & \textbf{$t$ (s)} & \textbf{gripper transition} & \textbf{$|F|$ (N)} \\
\midrule
gripper close (grasp)   & 16.74 & open $\to$ closed & 3.1 \\
force contact (press)   & 24.77 & closed (holding)  & 26.4 \\
gripper open (release)  & 29.01 & closed $\to$ open  & 1.9 \\
\bottomrule
\end{tabular}
\caption{Detected interaction events. The gripper-state transition is the
detection signal; force magnitude is reported at the event. Times are taken at
the threshold crossing (grasp/release) and at the force peak (press).}
\end{table}

\begin{figure}[H]
\centering
\includegraphics[width=\textwidth]{figures/timeline.png}
\caption{Proprioceptive timeline. Gripper width (blue) drops sharply at the
grasp ($\sim$16.7\,s) and rises at release ($\sim$29\,s); force magnitude (red)
peaks at $\sim$26\,N during the press onto the cup ($\sim$24.8\,s). The
pick/press/drop phases are unambiguous from the signals alone.}
\label{fig:timeline}
\end{figure}

\begin{figure}[H]
\centering
\includegraphics[width=\textwidth]{figures/event_frames.png}
\caption{ZED frames at the detected event timestamps. \textbf{Left:} grasp ---
the carrot is in the gripper above the plate. \textbf{Centre:} release --- the
carrot at the cup. \textbf{Right:} press --- the gripper forcing the carrot onto
the cup. Each event frame contains exactly the object the robot is interacting
with, with no vision processing yet applied.}
\label{fig:frames}
\end{figure}

This validates the core assumption before any segmentation model is introduced:
proprioceptive events are reliable temporal anchors, and the corresponding
frames localise the relevant object in time.

\section{Prospective Research Framing}

The lab's existing learning-from-demonstration pipeline forms a clear arc:
kinematic constraint identification from pose
(Overbeek et al., \emph{Frontiers in Robotics and AI}, 2025) and physical-interaction
identification from contact (Overbeek et al., \emph{Advanced Robotics Research},
2025). A natural third layer is \textbf{visual object grounding}: identifying not
just \emph{that} a constraint or contact occurs, but \emph{which object} it
concerns.

\paragraph{Proposed contribution.} Use gripper and force/contact events to
disambiguate which object in a label-free proposal set is task-relevant per
sub-phase --- combining proprioceptive event detection with a promptable
segmenter (SAM/SAM\,2, prompted by the end-effector projected into the image) or
a depth+attention proposer (DADO). Prior work uses gripper events alone to learn
affordances, or depth/attention alone to propose objects; the combination of
\emph{contact-force events} with label-free object discovery for per-phase
task-relevant grounding appears to be an open gap, and a well-scoped story for a
short paper (e.g.\ RA-L), with hardware validation on the lab's Franka setup.

\paragraph{Next steps.}
\begin{enumerate}[noitemsep]
  \item Run SAM/SAM\,2 (and compare DADO) on the three event frames; turn
        ``object visible'' into ``object segmented.''
  \item Project the end-effector pose into the camera frame as a spatial prompt
        (requires the ZED-to-base extrinsics).
  \item Propagate masks across the demo (SAM\,2 video) and check consistency
        against the proprioceptive timeline.
  \item Scale to multiple trials/tasks once more data is collected; quantify
        against gripper-only and vision-only baselines.
\end{enumerate}

\clearpage

\section{Update --- Segmentation, Propagation, and a First Deliverable}

This section extends the writeup with results obtained after the first version
was sent. The original four pages are unchanged. Steps 1 and 3 from the previous
section are now complete; step 2 remains blocked on calibration data from the
lab; step 4 is blocked on additional trials.

\begin{figure}[H]
\centering
\resizebox{\textwidth}{!}{% Pipeline diagram body (no \documentclass wrapper).
% Included from Docs/writeup.tex via % Pipeline diagram body (no \documentclass wrapper).
% Included from Docs/writeup.tex via % Pipeline diagram body (no \documentclass wrapper).
% Included from Docs/writeup.tex via \input{figures_tex/pipeline_inner.tex}
% and from Docs/figures_tex/pipeline.tex (standalone) for PDF/PNG export.

\definecolor{inFill}{RGB}{230,230,230}
\definecolor{inBorder}{RGB}{140,140,140}
\definecolor{coreFill}{RGB}{255,225,185}
\definecolor{coreBorder}{RGB}{210,130,40}
\definecolor{frozenFill}{RGB}{210,225,250}
\definecolor{frozenBorder}{RGB}{90,130,190}
\definecolor{outFill}{RGB}{215,240,215}
\definecolor{outBorder}{RGB}{60,140,60}
\definecolor{evColor}{RGB}{190,50,50}

\begin{tikzpicture}[
  inblock/.style={
    draw=inBorder, thick, rounded corners=3pt, fill=inFill,
    minimum width=2.4cm, minimum height=1.5cm, text width=2.2cm, align=center, font=\small
  },
  core/.style={
    draw=coreBorder, very thick, rounded corners=3pt, fill=coreFill,
    minimum width=2.6cm, minimum height=2.0cm, text width=2.4cm, align=center, font=\small
  },
  frozen/.style={
    draw=frozenBorder, very thick, rounded corners=3pt, dashed, fill=frozenFill,
    minimum width=2.8cm, minimum height=2.0cm, text width=2.6cm, align=center, font=\small
  },
  outblock/.style={
    draw=outBorder, very thick, rounded corners=3pt, fill=outFill,
    minimum width=3.0cm, minimum height=2.0cm, text width=2.8cm, align=center, font=\small
  },
  arr/.style={-{Stealth[length=2.5mm]}, thick},
  evarr/.style={-{Stealth[length=2.5mm]}, very thick, evColor},
  rlabel/.style={font=\scriptsize\sffamily, fill=white, inner sep=1.5pt},
  dlabel/.style={font=\scriptsize, align=center},
  node distance=1.4cm,
]

% --- Row 1: inputs ---
\node[inblock] (bag) {\textbf{ROS\,2 bag}\\[2pt]{\scriptsize pose / gripper /}\\{\scriptsize wrench / RGB}};
\node[inblock, right=1.3cm of bag] (csv) {\textbf{Merged CSV}\\[2pt]{\scriptsize pose-synced,}\\{\scriptsize topic-prefixed}};
\node[inblock, right=1.3cm of csv] (frames) {\textbf{ZED PNGs}\\[2pt]{\scriptsize 497 frames,}\\{\scriptsize ns-stamped}};

% --- Row 2: core stages, horizontal flow ---
\node[core, below=2.6cm of bag] (events) {\textbf{Event}\\\textbf{Detection}\\[3pt]
  {\scriptsize gripper-width thr.}\\
  {\scriptsize $\Rightarrow$ grasp, release}\\
  {\scriptsize force peak}\\
  {\scriptsize $\Rightarrow$ press}};

\node[core, right=1.5cm of events] (seed) {\textbf{Frame +}\\\textbf{Prompt Seed}\\[3pt]
  {\scriptsize event time $\to$}\\
  {\scriptsize image ts $\to$}\\
  {\scriptsize point prompt}};

\node[frozen, right=1.5cm of seed] (sam) {\textbf{SAM\,2 (frozen)}\\\textbf{Video Pred.}\\[3pt]
  {\scriptsize bidirectional}\\
  {\scriptsize propagation}\\
  {\scriptsize per object}};

\node[core, right=1.5cm of sam] (agg) {\textbf{Per-Frame}\\\textbf{Aggregation}\\[3pt]
  {\scriptsize obj 1: grasped}\\
  {\scriptsize obj 2: contact-rec.}\\
  {\scriptsize bbox + area + role}};

% --- Row 3: outputs ---
\node[outblock, below=2.3cm of seed] (json) {\textbf{Object Sidecar}\\[3pt]
  {\scriptsize \texttt{objects.json}}\\
  {\scriptsize \texttt{objects\_summary.csv}}\\
  {\scriptsize overlays + MP4}};

\node[inblock, right=2.5cm of json] (downstream) {\textbf{Downstream}\\\textbf{LfD Pipeline}\\[2pt]{\scriptsize constraint +}\\{\scriptsize interaction ID}};

% --- Arrows: row 1 -> row 2 ---
\draw[arr] (bag) -- node[above, rlabel] {\texttt{ros2\_unbag}} (csv);
\draw[arr] (csv) -- (frames);
\draw[arr] (csv.south) -- node[left, rlabel, pos=0.6] {pose, gripper, wrench} (events.north);
\draw[arr] (frames.south) |- ([yshift=0.5cm]seed.north) -- node[right, rlabel, pos=0.4] {image-ts col} (seed.north);

% --- Arrows: row 2 horizontal flow ---
\draw[evarr] (events) -- node[above, rlabel] {3 events} (seed);
\draw[arr] (seed) -- node[above, rlabel] {seed prompts} (sam);
\draw[arr] (sam) -- node[above, rlabel] {per-frame masks} (agg);

% --- Arrows: row 2 -> row 3 ---
\draw[arr] (agg.south) |- ([yshift=0.4cm]json.east) -- (json.east);
\draw[arr] (json) -- node[above, rlabel] {consumed by} (downstream);

\begin{pgfonlayer}{background}
  \node[draw=inBorder, dashed, thick, rounded corners=5pt, fill=gray!4, inner sep=10pt,
        fit=(bag)(csv)(frames),
        label={[font=\scriptsize\sffamily, text=inBorder]above:Inputs (from lab export)}
  ] {};
  \node[draw=coreBorder, thick, rounded corners=5pt, fill=orange!4, inner sep=10pt,
        fit=(events)(seed)(sam)(agg),
        label={[font=\scriptsize\sffamily\bfseries, text=coreBorder]above:Proprioception-cued logic (this work)}
  ] {};
\end{pgfonlayer}

\end{tikzpicture}

% and from Docs/figures_tex/pipeline.tex (standalone) for PDF/PNG export.

\definecolor{inFill}{RGB}{230,230,230}
\definecolor{inBorder}{RGB}{140,140,140}
\definecolor{coreFill}{RGB}{255,225,185}
\definecolor{coreBorder}{RGB}{210,130,40}
\definecolor{frozenFill}{RGB}{210,225,250}
\definecolor{frozenBorder}{RGB}{90,130,190}
\definecolor{outFill}{RGB}{215,240,215}
\definecolor{outBorder}{RGB}{60,140,60}
\definecolor{evColor}{RGB}{190,50,50}

\begin{tikzpicture}[
  inblock/.style={
    draw=inBorder, thick, rounded corners=3pt, fill=inFill,
    minimum width=2.4cm, minimum height=1.5cm, text width=2.2cm, align=center, font=\small
  },
  core/.style={
    draw=coreBorder, very thick, rounded corners=3pt, fill=coreFill,
    minimum width=2.6cm, minimum height=2.0cm, text width=2.4cm, align=center, font=\small
  },
  frozen/.style={
    draw=frozenBorder, very thick, rounded corners=3pt, dashed, fill=frozenFill,
    minimum width=2.8cm, minimum height=2.0cm, text width=2.6cm, align=center, font=\small
  },
  outblock/.style={
    draw=outBorder, very thick, rounded corners=3pt, fill=outFill,
    minimum width=3.0cm, minimum height=2.0cm, text width=2.8cm, align=center, font=\small
  },
  arr/.style={-{Stealth[length=2.5mm]}, thick},
  evarr/.style={-{Stealth[length=2.5mm]}, very thick, evColor},
  rlabel/.style={font=\scriptsize\sffamily, fill=white, inner sep=1.5pt},
  dlabel/.style={font=\scriptsize, align=center},
  node distance=1.4cm,
]

% --- Row 1: inputs ---
\node[inblock] (bag) {\textbf{ROS\,2 bag}\\[2pt]{\scriptsize pose / gripper /}\\{\scriptsize wrench / RGB}};
\node[inblock, right=1.3cm of bag] (csv) {\textbf{Merged CSV}\\[2pt]{\scriptsize pose-synced,}\\{\scriptsize topic-prefixed}};
\node[inblock, right=1.3cm of csv] (frames) {\textbf{ZED PNGs}\\[2pt]{\scriptsize 497 frames,}\\{\scriptsize ns-stamped}};

% --- Row 2: core stages, horizontal flow ---
\node[core, below=2.6cm of bag] (events) {\textbf{Event}\\\textbf{Detection}\\[3pt]
  {\scriptsize gripper-width thr.}\\
  {\scriptsize $\Rightarrow$ grasp, release}\\
  {\scriptsize force peak}\\
  {\scriptsize $\Rightarrow$ press}};

\node[core, right=1.5cm of events] (seed) {\textbf{Frame +}\\\textbf{Prompt Seed}\\[3pt]
  {\scriptsize event time $\to$}\\
  {\scriptsize image ts $\to$}\\
  {\scriptsize point prompt}};

\node[frozen, right=1.5cm of seed] (sam) {\textbf{SAM\,2 (frozen)}\\\textbf{Video Pred.}\\[3pt]
  {\scriptsize bidirectional}\\
  {\scriptsize propagation}\\
  {\scriptsize per object}};

\node[core, right=1.5cm of sam] (agg) {\textbf{Per-Frame}\\\textbf{Aggregation}\\[3pt]
  {\scriptsize obj 1: grasped}\\
  {\scriptsize obj 2: contact-rec.}\\
  {\scriptsize bbox + area + role}};

% --- Row 3: outputs ---
\node[outblock, below=2.3cm of seed] (json) {\textbf{Object Sidecar}\\[3pt]
  {\scriptsize \texttt{objects.json}}\\
  {\scriptsize \texttt{objects\_summary.csv}}\\
  {\scriptsize overlays + MP4}};

\node[inblock, right=2.5cm of json] (downstream) {\textbf{Downstream}\\\textbf{LfD Pipeline}\\[2pt]{\scriptsize constraint +}\\{\scriptsize interaction ID}};

% --- Arrows: row 1 -> row 2 ---
\draw[arr] (bag) -- node[above, rlabel] {\texttt{ros2\_unbag}} (csv);
\draw[arr] (csv) -- (frames);
\draw[arr] (csv.south) -- node[left, rlabel, pos=0.6] {pose, gripper, wrench} (events.north);
\draw[arr] (frames.south) |- ([yshift=0.5cm]seed.north) -- node[right, rlabel, pos=0.4] {image-ts col} (seed.north);

% --- Arrows: row 2 horizontal flow ---
\draw[evarr] (events) -- node[above, rlabel] {3 events} (seed);
\draw[arr] (seed) -- node[above, rlabel] {seed prompts} (sam);
\draw[arr] (sam) -- node[above, rlabel] {per-frame masks} (agg);

% --- Arrows: row 2 -> row 3 ---
\draw[arr] (agg.south) |- ([yshift=0.4cm]json.east) -- (json.east);
\draw[arr] (json) -- node[above, rlabel] {consumed by} (downstream);

\begin{pgfonlayer}{background}
  \node[draw=inBorder, dashed, thick, rounded corners=5pt, fill=gray!4, inner sep=10pt,
        fit=(bag)(csv)(frames),
        label={[font=\scriptsize\sffamily, text=inBorder]above:Inputs (from lab export)}
  ] {};
  \node[draw=coreBorder, thick, rounded corners=5pt, fill=orange!4, inner sep=10pt,
        fit=(events)(seed)(sam)(agg),
        label={[font=\scriptsize\sffamily\bfseries, text=coreBorder]above:Proprioception-cued logic (this work)}
  ] {};
\end{pgfonlayer}

\end{tikzpicture}

% and from Docs/figures_tex/pipeline.tex (standalone) for PDF/PNG export.

\definecolor{inFill}{RGB}{230,230,230}
\definecolor{inBorder}{RGB}{140,140,140}
\definecolor{coreFill}{RGB}{255,225,185}
\definecolor{coreBorder}{RGB}{210,130,40}
\definecolor{frozenFill}{RGB}{210,225,250}
\definecolor{frozenBorder}{RGB}{90,130,190}
\definecolor{outFill}{RGB}{215,240,215}
\definecolor{outBorder}{RGB}{60,140,60}
\definecolor{evColor}{RGB}{190,50,50}

\begin{tikzpicture}[
  inblock/.style={
    draw=inBorder, thick, rounded corners=3pt, fill=inFill,
    minimum width=2.4cm, minimum height=1.5cm, text width=2.2cm, align=center, font=\small
  },
  core/.style={
    draw=coreBorder, very thick, rounded corners=3pt, fill=coreFill,
    minimum width=2.6cm, minimum height=2.0cm, text width=2.4cm, align=center, font=\small
  },
  frozen/.style={
    draw=frozenBorder, very thick, rounded corners=3pt, dashed, fill=frozenFill,
    minimum width=2.8cm, minimum height=2.0cm, text width=2.6cm, align=center, font=\small
  },
  outblock/.style={
    draw=outBorder, very thick, rounded corners=3pt, fill=outFill,
    minimum width=3.0cm, minimum height=2.0cm, text width=2.8cm, align=center, font=\small
  },
  arr/.style={-{Stealth[length=2.5mm]}, thick},
  evarr/.style={-{Stealth[length=2.5mm]}, very thick, evColor},
  rlabel/.style={font=\scriptsize\sffamily, fill=white, inner sep=1.5pt},
  dlabel/.style={font=\scriptsize, align=center},
  node distance=1.4cm,
]

% --- Row 1: inputs ---
\node[inblock] (bag) {\textbf{ROS\,2 bag}\\[2pt]{\scriptsize pose / gripper /}\\{\scriptsize wrench / RGB}};
\node[inblock, right=1.3cm of bag] (csv) {\textbf{Merged CSV}\\[2pt]{\scriptsize pose-synced,}\\{\scriptsize topic-prefixed}};
\node[inblock, right=1.3cm of csv] (frames) {\textbf{ZED PNGs}\\[2pt]{\scriptsize 497 frames,}\\{\scriptsize ns-stamped}};

% --- Row 2: core stages, horizontal flow ---
\node[core, below=2.6cm of bag] (events) {\textbf{Event}\\\textbf{Detection}\\[3pt]
  {\scriptsize gripper-width thr.}\\
  {\scriptsize $\Rightarrow$ grasp, release}\\
  {\scriptsize force peak}\\
  {\scriptsize $\Rightarrow$ press}};

\node[core, right=1.5cm of events] (seed) {\textbf{Frame +}\\\textbf{Prompt Seed}\\[3pt]
  {\scriptsize event time $\to$}\\
  {\scriptsize image ts $\to$}\\
  {\scriptsize point prompt}};

\node[frozen, right=1.5cm of seed] (sam) {\textbf{SAM\,2 (frozen)}\\\textbf{Video Pred.}\\[3pt]
  {\scriptsize bidirectional}\\
  {\scriptsize propagation}\\
  {\scriptsize per object}};

\node[core, right=1.5cm of sam] (agg) {\textbf{Per-Frame}\\\textbf{Aggregation}\\[3pt]
  {\scriptsize obj 1: grasped}\\
  {\scriptsize obj 2: contact-rec.}\\
  {\scriptsize bbox + area + role}};

% --- Row 3: outputs ---
\node[outblock, below=2.3cm of seed] (json) {\textbf{Object Sidecar}\\[3pt]
  {\scriptsize \texttt{objects.json}}\\
  {\scriptsize \texttt{objects\_summary.csv}}\\
  {\scriptsize overlays + MP4}};

\node[inblock, right=2.5cm of json] (downstream) {\textbf{Downstream}\\\textbf{LfD Pipeline}\\[2pt]{\scriptsize constraint +}\\{\scriptsize interaction ID}};

% --- Arrows: row 1 -> row 2 ---
\draw[arr] (bag) -- node[above, rlabel] {\texttt{ros2\_unbag}} (csv);
\draw[arr] (csv) -- (frames);
\draw[arr] (csv.south) -- node[left, rlabel, pos=0.6] {pose, gripper, wrench} (events.north);
\draw[arr] (frames.south) |- ([yshift=0.5cm]seed.north) -- node[right, rlabel, pos=0.4] {image-ts col} (seed.north);

% --- Arrows: row 2 horizontal flow ---
\draw[evarr] (events) -- node[above, rlabel] {3 events} (seed);
\draw[arr] (seed) -- node[above, rlabel] {seed prompts} (sam);
\draw[arr] (sam) -- node[above, rlabel] {per-frame masks} (agg);

% --- Arrows: row 2 -> row 3 ---
\draw[arr] (agg.south) |- ([yshift=0.4cm]json.east) -- (json.east);
\draw[arr] (json) -- node[above, rlabel] {consumed by} (downstream);

\begin{pgfonlayer}{background}
  \node[draw=inBorder, dashed, thick, rounded corners=5pt, fill=gray!4, inner sep=10pt,
        fit=(bag)(csv)(frames),
        label={[font=\scriptsize\sffamily, text=inBorder]above:Inputs (from lab export)}
  ] {};
  \node[draw=coreBorder, thick, rounded corners=5pt, fill=orange!4, inner sep=10pt,
        fit=(events)(seed)(sam)(agg),
        label={[font=\scriptsize\sffamily\bfseries, text=coreBorder]above:Proprioception-cued logic (this work)}
  ] {};
\end{pgfonlayer}

\end{tikzpicture}
}
\caption{Pipeline overview. The lab's standard ROS\,2 bag export (grey,
top) is consumed as-is. Proprioceptive event detection on the merged CSV
identifies the grasp, press, and release moments; each event indexes the
corresponding ZED frame, which is used to seed SAM\,2 (frozen) with a point
prompt. Bidirectional video propagation yields per-frame, role-tagged masks
that are aggregated into a JSON sidecar consumable by downstream LfD code.}
\label{fig:pipeline}
\end{figure}

\subsection{Step 1 (revisited): SAM and SAM\,2 on the event frames}

Running \textbf{SAM (ViT-H)} on the three event frames with a point prompt
placed near the gripper produced clean segmentations once the prompt sits on
the relevant object (Figure~\ref{fig:sam}). The grasp frame yields a tight
carrot mask (predictor score $1.01$); the press and release frames yield cup
masks of varying completeness depending on point placement. The takeaway is
unsurprising but important: SAM is a competent paintbrush; the open question is
where to point it.

\begin{figure}[H]
\centering
\includegraphics[width=\textwidth]{figures/segmented_events.png}
\caption{SAM (ViT-H) with point prompts on the three event frames. Left: the
carrot is recovered cleanly at the grasp event. Centre and right: the cup is
partially recovered at the press and release events. The blue dot marks the
prompt point in each frame.}
\label{fig:sam}
\end{figure}

\subsection{Step 3: SAM\,2 mask propagation across the demo}

Seeding \textbf{SAM\,2 (Hiera-L)} with a single point prompt on the carrot at
the grasp event and propagating forward through the remaining 247 frames
recovers the carrot mask across the entire pick~$\rightarrow$~press~$\rightarrow$~drop
sequence. A second SAM\,2 state, seeded on the cup at the press event and
propagated bidirectionally, recovers the cup mask across all 497 frames.
Figure~\ref{fig:both_strip} shows six milestone frames with both objects
overlaid.

\begin{figure}[H]
\centering
\includegraphics[width=\textwidth]{figures/propagation_both_strip.png}
\caption{Per-frame masks for the grasped object (carrot, green) and
contact-receiving object (cup, magenta), shown at six milestones spanning the
demo. Both masks come from SAM\,2 video propagation seeded at the corresponding
proprioceptive event (grasp for the carrot, press for the cup). No mask is
hand-labelled.}
\label{fig:both_strip}
\end{figure}

\subsection{Quantitative read --- mask area over time}

Mask area is a useful summary signal. Figure~\ref{fig:area} plots both objects
on the relative-time axis with the three events marked. Two observations:

\begin{itemize}[noitemsep]
  \item The cup mask area peaks essentially at the press event, dropping
        sharply once the carrot is released onto it.
  \item The carrot mask shrinks from $\sim 38{,}000$\,px at grasp to a stable
        $\sim 20{,}000$\,px once it is in the gripper (occlusion by the
        gripper itself), and collapses after release. Absolute area is
        therefore a coarse ``presence'' indicator; bounding boxes or
        centroids are the more useful per-frame summary.
\end{itemize}

\begin{figure}[H]
\centering
\includegraphics[width=\textwidth]{figures/mask_area_over_time.png}
\caption{Per-frame mask area for the carrot (green) and the cup (magenta),
aligned to the proprioceptive event times. The cup's visual prominence
co-rises with the force-contact event.}
\label{fig:area}
\end{figure}

\subsection{Comparison --- a label-free baseline (DADO-style)}

For comparison, we ran a DADO-style label-free object proposer on the three
event frames: DINOv2 self-attention combined with monocular depth
(Depth-Anything-V2). The depth map cleanly separates foreground from
background, but the DINOv2 attention map is uniformly low and noisy on this
in-lab scene, so the attention $\times$ depth saliency does not isolate either
object (Figure~\ref{fig:dado}). The proposal masks are essentially empty.

\begin{figure}[H]
\centering
\includegraphics[width=\textwidth]{figures/dado_events.png}
\caption{DADO-style proposer on the three event frames (rows: grasp, press,
release; columns: RGB, DINOv2 attention, Depth-Anything-V2 depth, thresholded
saliency mask). Depth is informative; attention is not. The proposal masks
fail to isolate either the carrot or the cup. This validates that the
proprioceptive cue is load-bearing rather than decorative: vision-only
label-free proposal does not solve the task on this scene.}
\label{fig:dado}
\end{figure}

\subsection{Force-direction sanity check (uncalibrated)}

A small experiment toward step 2 above. At the press event the wrist wrench
is $|F| \approx 11.2$\,N, oriented essentially downward in the robot base
frame ($f_z = -11.1$\,N, lateral components within $\pm 1$\,N). To preview
what the end-effector projection will look like once real ZED-to-base
extrinsics are available, we fit a coarse linear base-to-image map
($\mathbb{R}^4 \rightarrow \mathbb{R}^2$) from the carrot-mask centroids over
time and use it to draw the force-reaction arrow on the press frame
(Figure~\ref{fig:force}). With the latest (auto-seeded) propagation,
the fit improved from RMSE $58$\,px (hand-seeded run) to $23$\,px over
$240$ paired points --- still uncalibrated, but consistent enough that the
projected arrow now stays within the press region. Real intrinsics +
extrinsics would replace this approximation with a geometric projection.

\begin{figure}[H]
\centering
\includegraphics[width=0.7\textwidth]{figures/force_overlay_press.png}
\caption{Press frame with carrot mask (green), cup mask (magenta), and the
force-reaction arrow (red), drawn through an uncalibrated linear base-to-image
fit. Fit RMSE = $23$\,px over $240$ carrot-centroid points (down from $58$\,px
with the original hand-seeded propagation). Real ZED-to-base extrinsics would
replace this approximation with a geometric projection.}
\label{fig:force}
\end{figure}

\subsection{Deliverable --- \texttt{identify\_objects.py} and the JSON sidecar}

The pipeline above is wrapped as a script intended to run on any future bag
produced by the lab's standard export. The output is a per-trial bundle
consumable by downstream LfD code:

\begin{itemize}[noitemsep]
  \item \texttt{objects.json} --- top-level events (grasp / press / release)
        with their relative times, row indices, and image timestamps; per-frame
        entries listing the object IDs present, their roles
        (\texttt{grasped} / \texttt{contact\_receiver}), pixel mask area, and
        axis-aligned bounding box.
  \item \texttt{objects\_summary.csv} --- the same content as a flat,
        spreadsheet-readable table (744 rows for this trial).
  \item \texttt{overlays/f0000\_*.png} \ldots --- per-frame overlay PNGs
        keyed by frame index.
  \item \texttt{overlay.mp4} --- a stitched 15\,Hz video for visual review.
\end{itemize}

An example entry from \texttt{objects.json} at the press frame
($\texttt{frame\_idx}=370$, $t_{\text{rel}}=24.76$\,s):

\begin{table}[H]
\centering
\begin{tabular}{lllr}
\toprule
\textbf{obj\_id} & \textbf{role} & \textbf{bbox (xyxy)} & \textbf{mask\,px} \\
\midrule
1 & grasped          & $[10,\,0,\,704,\,163]$  & 23{,}123 \\
2 & contact\_receiver & $[10,\,16,\,721,\,442]$ & 65{,}601 \\
\bottomrule
\end{tabular}
\caption{Per-frame object entry written by \texttt{identify\_objects.py} (press
event, frame 370). The JSON sidecar contains 497 such entries per trial.}
\end{table}

\subsection{Honest limitations}

Recorded here so they are visible to downstream readers; a more detailed list
is kept in \texttt{FAILURE\_MODES.md} in the repository.

\begin{itemize}[noitemsep]
  \item \textbf{Seed points: hand-fractioned defaults, auto-seed fallback.}
        The image-relative defaults in the propagation scripts (e.g.\
        $(0.70, 0.30)$ for the grasped object) are tuned to this trial.
        When \texttt{Code/auto\_seed.py} produces
        \texttt{figures/identify/\allowbreak auto\_seeds.csv}, both
        propagation scripts now consume it automatically; otherwise they
        fall back to the hard-coded defaults. The current auto-seed is a
        vision-only heuristic and is the source of the carrot drift below.
        Both routes are placeholders for end-effector projection, which
        requires ZED intrinsics + ZED-to-base extrinsics + the gripper/tool
        frame from the URDF.
  \item \textbf{Auto-seed drift on the grasped object.} On this trial the
        auto-seed for the grasped role landed on the gripper finger
        ($(514, 54)$ in a $960\times540$ frame), not on the carrot body, so
        the early propagation locks onto the gripper and the carrot mask
        only stabilises once the gripper moves. The cup auto-seed is clean.
        End-effector projection will replace the vision-only heuristic.
  \item \textbf{Carrot is not tracked during reach.} Forward propagation
        starts at the grasp; frames before grasp have no carrot mask. The
        cup is propagated bidirectionally so it is masked throughout.
  \item \textbf{DADO-style proposal is essentially empty here.} Pretrained
        DINOv2 attention does not isolate either object in this scene; a
        depth-aware proposer that uses depth as the primary cue (rather than
        as a re-weighting of attention) would likely fare better and is left
        for follow-up once depth is enabled on the bag.
  \item \textbf{One trial, one task, one camera pose.} All quantitative
        statements are evaluative on \texttt{lfdws\_t001} alone. Step 4
        of the original roadmap (multi-trial baselines) remains blocked on
        additional bags.
\end{itemize}

\subsection{Bag-agnostic utilities (added since the last update)}

Four small additions to make the pipeline run on future bags without
hand-tuning per trial:

\begin{itemize}[noitemsep]
  \item \texttt{Code/auto\_seed.py} --- replaces the hard-coded image
        fractions used to seed SAM\,2. Runs SAM's automatic mask generator
        on each event frame, scores candidates by centrality / size /
        border-distance and a small role prior (\textit{grasped} prefers
        upper-centre, \textit{contact\_receiver} prefers lower-centre),
        and writes the chosen seed point to
        \texttt{figures/identify/auto\_seeds.csv}. Both
        \texttt{Code/propagate\_demo.py} and \texttt{Code/propagate\_cup.py}
        now read this CSV automatically when it exists, falling back to the
        hard-coded defaults otherwise --- so a fresh bag goes through the
        pipeline with no per-trial constants to edit. On the carrot trial
        it recovers a clean cup seed and a gripper-finger seed near the
        carrot (see Honest Limitations).
  \item \texttt{Code/multi\_event.py} --- generalises event detection to
        find \emph{all} grasps, releases, and force-magnitude peaks above a
        baseline, then groups them into interaction cycles. On
        \texttt{lfdws\_t001} it returns one cycle with one grasp,
        one release, and four force peaks (the original detector returned
        only the single global peak). This is what unblocks bags with
        multi-cycle tasks.
  \item \texttt{Code/object\_identity.py} --- embeds per-frame mask crops
        with frozen DINOv2 and clusters the embeddings to assign stable
        object identities across phases (and, in future, across episodes).
        Crops are now tightened to the actual mask: the colour-coded
        overlay PNG is parsed back into a boolean mask, the crop is taken
        at the mask's bounding box, and the background inside the crop is
        blacked out so the embedding sees the object rather than the table.
        On this trial, $25$ sample crops at distance threshold $0.6$ yield
        $3$ stable identities (the carrot+gripper ensemble, a partial-cup
        cluster, and a partial-carrot cluster). Closer to the true count
        of two; remaining over-merging is dominated by the gripper-finger
        drift noted in Honest Limitations.
  \item \texttt{Code/trial\_report.py} --- builds a single PDF per trial
        (\texttt{Docs/reports/report\_<trial>.pdf}) containing the summary
        table, the event list, the press-frame sidecar excerpt, and every
        figure produced by the pipeline. Intended as the one document to
        share per future bag.
  \item \texttt{Code/project\_ee.py} --- the calibrated successor to the
        vision-only auto-seed and the uncalibrated force overlay. Given
        camera intrinsics, the \texttt{base}$\to$\texttt{camera} extrinsic,
        and the F/T-sensor mount (all read from \texttt{calibration.yaml}),
        it projects the end-effector point into the image as a geometric
        SAM seed and projects the wrench line of action (Bicchi's contact
        ray, $r_0 = (\mathbf{f}\times\boldsymbol{\tau})/|\mathbf{f}|^2$) to
        localise the contact-receiving object. It runs in a DRY mode that
        reports the base-frame geometry but draws nothing until the real
        calibration is supplied, so the pipeline never uses placeholder
        numbers. The Franka Research~3 arm URDF is now vendored
        (\texttt{Data/fr3.urdf}); \texttt{current\_pose} is published in the
        \texttt{base} frame and is the fingertip TCP under the default
        configuration.
\end{itemize}

\subsection{Prospective research direction (updated)}

The original framing on page 4 proposed combining contact-force events
with label-free object discovery as the contribution; after building the
pipeline end-to-end on \texttt{lfdws\_t001}, the picture has sharpened in
two ways worth recording before the next sync.

\paragraph{What the data has shown so far.}
\begin{itemize}[noitemsep]
  \item \textbf{The role split is empirically real, not just a framing
        trick.} Proprioceptive signatures cleanly separate the grasped
        object (gripper-state transition) from the contact-receiving
        object (wrench peak during a held window). The press signal does
        not coincide with any gripper transition --- this is what makes
        the two roles structurally distinguishable from a single demo,
        and it is the part the existing self-supervised literature
        (gripper-events only) does not address.
  \item \textbf{Vision-only label-free proposers fail on this scene.}
        DINOv2-attention $\times$ Depth-Anything saliency does not isolate
        either object (Section~5.4). The proprioceptive cue is therefore
        load-bearing, not decorative --- the pipeline is not just SAM
        with extra steps.
  \item \textbf{The wrench vector carries usable spatial information.}
        The uncalibrated base-to-image fit (Section~5.5) is suggestive
        rather than geometric, but it is concentrated enough to keep the
        force-reaction arrow inside the press region. With real
        intrinsics + extrinsics it becomes a wrench-projected segmentation
        prompt --- a distinct, methodological lever the prior art does
        not exploit.
\end{itemize}

\paragraph{Refined contribution claim.} A label-free system that, from a
single demonstration, automatically assigns \emph{role-tagged} per-frame
object masks --- \textit{grasped} from the gripper-state transition,
\textit{contact-receiving} from the wrench-magnitude peak --- where the
two roles correspond to two structurally distinct objects in the scene
without any cross-frame appearance learning, language input, or per-trial
hand tuning. Hardware validation on the lab's Franka. The intended
extension once camera calibration arrives: use the wrench \emph{direction}
(not just the event time) as a geometric prompt for promptable
segmentation, distinguishing the long-tool case (carrot, peg, stick)
from EE-only projection methods that assume the contact point is at the
gripper origin.

\paragraph{What stays open and is genuinely blocked on data.}
Generalisation across tasks, objects, and camera poses; a downstream
result showing that role-tagged masks improve the lab's existing
constraint-identification and interaction-identification pipelines; and
the wrench-projection variant above. None of these are answerable on
\texttt{lfdws\_t001} alone --- they are the reason for the next-sync asks
below.

\subsection{Asks for the next sync}

Concrete data items needed to remove the placeholders above:

\begin{enumerate}[noitemsep]
  \item ZED intrinsics (camera matrix $K$, distortion) and the
        \texttt{base}$\to$\texttt{camera} extrinsic transform --- fills
        \texttt{calibration.yaml} and takes \texttt{project\_ee.py} live.
  \item The F/T-sensor mount transform
        (\texttt{bota\_frame}$\to$\texttt{base}). This is the one quantity
        not recoverable from any URDF, and it is required for the
        wrench-line projection.
  \item Confirmation of whether \texttt{current\_pose} is the fingertip TCP
        or the flange (the script supports either).
  \item Depth stream enabled in the next bag.
  \item A re-centred camera, addressing the off-centre framing in the
        current trial.
  \item Two to three additional trials with varied objects and placements,
        to begin step 4.
\end{enumerate}

\noindent The Franka Research~3 arm URDF has since been obtained, so the
gripper/tool-frame request from the previous version is resolved.

\clearpage

\section{Update --- A Second Trial, and What It Revealed}

This section covers work since the last update, on a second recording,
\texttt{lfdws\_t001\_depth}. Pages 1--4 remain unchanged; the previous
update section (\S5) is also left as-is.

\subsection{Extracting a bag the standard exporter couldn't}

A depth-enabled recording arrived, but \texttt{ros2\_unbag} could not export
it: the ZED \texttt{compressedDepth} topic uses a transport it doesn't
support. Rather than wait on a lab-side plugin fix, we wrote
\texttt{Code/mcap\_extract.py}: a standalone \texttt{.mcap}$\to$merged-CSV
extractor requiring no ROS\,2 installation (\texttt{.mcap} files embed their
own message schemas, so the pure-Python \texttt{mcap-ros2-support} package
decodes \texttt{WrenchStamped}/\texttt{PoseStamped}/\texttt{CompressedImage}
directly). The \texttt{compressedDepth} transport itself was decoded by hand
(a 12-byte header of three \texttt{float32}s --- format flag, two
quantization constants --- followed by a PNG; true depth recovered as
$\text{depth}_m = \text{quantA}/(\text{code} - \text{quantB})$, written out
as 16-bit millimetre PNGs). Result: 1013/1013 RGB frames and 1013/1013 depth
frames decoded cleanly, output in the same layout the rest of the pipeline
already expects, so nothing downstream needed to change.

\subsection{The second trial is a different, richer task}

\texttt{lfdws\_t001\_depth} turned out not to be a carrot-pick-place at all.
The scene contains a Rubik's cube, a large plate, a black charger, and an
orange-handled screwdriver clamped to the table --- and the bag has no
gripper topic, so grasp/release events cannot be detected at all (only
force-magnitude peaks). A force-only event detector
(\texttt{Code/force\_only\_events.py}) found 16 peaks between $t{=}30$--$64$s;
grouping them by time-gap
(\texttt{Code/force\_only\_multi\_event.py}) resolved 5 distinct activity
clusters. Visually inspecting the seed frame of each cluster showed a
genuine multi-object, multi-step task: a press on the plate ($|F|{=}64$\,N,
the strongest contact), contact with the clamped screwdriver, a grasp of the
charger, and what looks like the charger being moved toward the screwdriver
clamp (a docking/plugging action). One of the five clusters
($t{\approx}44$--$47$s, contact with the screwdriver) falls exactly where the
plate's tracked mask area reads near-zero --- proprioception caught a real
interaction that the vision pipeline, seeded only on the plate, had no way
to see.

This matters beyond this one trial: it is direct evidence that the
pipeline's current two-role model (one \textit{grasped} object, one
\textit{contact\_receiver}, tied to exactly one grasp event and one press
event) does not generalise to tasks with more than two task-relevant objects
or more than one contact phase --- which, on this evidence, may be the norm
rather than the exception once more bags arrive.

\subsection{A concrete case of the seed-picking failure mode we'd flagged}

The auto-seed heuristic (\texttt{Code/auto\_seed.py}) scores SAM's candidate
masks by position and size, with a hard cutoff rejecting anything covering
more than 40\% of the frame --- a constant tuned to the first trial's small
cup. On this trial the correct contact-receiver (the plate) is large enough
to be rejected by that cutoff outright, and the picker fell back to a
meaningless sliver of background table, which then propagated through all
1013 frames as a wildly unstable mask (17--78\% of the frame, tracking
reflective table texture rather than any object). We manually reseeded the
plate and reran propagation; the corrected result is physically sane (peaks
at 31\% of frame right at the press, decays cleanly through occlusion) and
was verified against the raw frames, not just the summary numbers.

Two independent attempts at a general, non-manual fix were tried and both
failed, which is itself informative. A relative-size scoring scheme (reward
the largest candidate mask in the frame, rather than an absolute area cutoff)
fixes this trial but breaks the first one: on \texttt{lfdws\_t001}, the
reflective table mat is simply the largest candidate in that scene too, so
the same heuristic mis-picks it as the cup. A depth-based near-camera prior
(using this trial's real depth data, unavailable for calibration purposes but
usable for relative comparisons) also failed, because ZED depth dropout is
worst on exactly the specular surfaces causing the position/size confusion in
the first place --- the one signal that could rescue the pick is unavailable
precisely where it's needed. Both attempts were reverted; production
behaviour is unchanged. This strengthens rather than weakens the case that
end-effector projection (\texttt{Code/project\_ee.py}, still blocked on
calibration) is the correct fix, not a SAM-heuristic tuning problem.

\subsection{Generalising to N objects, demonstrated not just designed}

To show the two-role limitation above is fixable, we wrote a generic,
parameterised propagation script (\texttt{Code/propagate\_object\_n.py} ---
takes \texttt{obj\_id}/\texttt{role}/seed point as arguments, not hardcoded
to a fixed object) and a sidecar builder that accepts an arbitrary list of
objects instead of a fixed pair
(\texttt{Code/build\_sidecar\_multi.py}). The screwdriver and the charger
were both manually seeded and propagated bidirectionally across all 1013
frames (mask areas 0--43{,}000\,px and 0--46{,}700\,px respectively, no
runaway on either), and composed into a genuine three-object sidecar
alongside the corrected plate mask (Figure~\ref{fig:multi}) --- all three
objects visible and correctly masked simultaneously at the frame where
their extents overlap most. That covers 3 of the $\geq\!4$ objects
identified in \S6.2 (the docking phase, where the charger is moved toward
the screwdriver rather than being a distinct new object, was not
separately propagated); the schema and scripts support any $N$.

\begin{figure}[H]
\centering
\includegraphics[width=0.75\textwidth]{figures/identify_depth_multi/overlays/f0786_1782835535683396106.png}
\caption{Three independently-seeded, independently-propagated objects
(plate, magenta; screwdriver, orange; charger, yellow) composed into one
sidecar frame by \texttt{Code/build\_sidecar\_multi.py}. Proof that the
pipeline generalises past the original grasped/contact-receiver pair.}
\label{fig:multi}
\end{figure}

\subsection{Other results from this trial}

\begin{itemize}[noitemsep]
  \item \textbf{Real depth vs.\ estimated depth for label-free proposal.}
        Re-running the DADO-style comparison (\S5.4) with this trial's real
        sensor depth in place of Depth-Anything's estimate, across all 5
        activity-cluster frames (\texttt{Code/\_dado\_real\_depth\_multi.py}):
        mixed result. ZED depth validity itself varies sharply by surface
        (16\% valid on the specular plate vs.\ 58--72\% on the matte
        tool/charger), and the resulting object proposal genuinely lands on
        the correct object for 2 of 5 events, mostly hitting background
        scatter for the other 3. This is a more nuanced negative result than
        the single-frame version suggested, and it still supports the core
        claim: proprioceptive cueing remains necessary even with real depth.
  \item \textbf{Object identity does not generalise across trials, at
        least not yet.} Pooling DINOv2 crop embeddings from both trials
        (\texttt{Code/object\_identity\_cross\_trial.py}) into one clustering
        run fragments the 3 true objects into 12 clusters and mixes 2 of
        them together. A follow-up sweep across crop styles and thresholds
        (\texttt{Code/object\_identity\_cross\_trial\_sweep.py}) found that
        keeping background context (rather than the tight, blacked-out crops
        used previously) eliminates cross-object confusion entirely across a
        range of thresholds, but does not fix fragmentation --- the same
        object's crops across pose/time still don't cluster together. The
        two problems have different causes and need different fixes.
  \item \textbf{Bidirectional propagation and a presence signal, applied
        to the first trial too.} The grasped-object propagation on
        \texttt{lfdws\_t001} now runs backward+forward from the grasp
        event like the cup already did, giving full-demo coverage instead
        of forward-only. A bounding-box-diagonal presence signal was also
        added alongside raw mask area, and is visibly more stable under
        partial occlusion.
\end{itemize}

\subsection{Updated asks}

The calibration ask from \S5.9 stands (ZED intrinsics, \texttt{base}
$\to$\texttt{camera} extrinsics, and the \texttt{bota\_frame}$\to$\texttt{base}
F/T mount transform) --- everything in this section that touches seed
placement or the wrench projection remains blocked on it. Two new,
trial-specific questions: whether the missing gripper topic on
\texttt{lfdws\_t001\_depth} was intentional, and --- given \S6.2's finding
--- whether future recordings can be scoped to fewer simultaneous
objects/phases per bag, since the current pipeline's role model handles
$N{=}3$ but has not yet been exercised at the scale this trial actually
needs.

\clearpage

\section{Update --- Two More Trials, a Third Task}

Two new recordings arrived, \texttt{lfdws\_t004} and \texttt{lfdws\_t005},
extracted with the same \texttt{Code/mcap\_extract.py} pipeline as before.
Both carry the gripper topic (unlike \texttt{lfdws\_t001\_depth}), so real
grasp/release detection applies directly, but neither has any wrench data
at all --- the F/T sensor published zero messages on both --- so no
\textit{press} event is detectable on either trial. Event detection across
the pipeline was extended to handle this gripper-only case symmetrically
to the existing force-only case, and regression-tested against both prior
trials before trusting it on the new ones.

\subsection{A third distinct task}

Neither trial is the carrot/cup or plate/screwdriver scene from before.
Both show a small fastener or knob being picked up and inserted into a
mounting hole on a wall-mounted pegboard, near a hinge/latch assembly ---
the same lab rig as the depth trial, reused for a different manipulation.

\begin{table}[H]
\centering
\begin{tabular}{lllll}
\toprule
\textbf{Trial} & \textbf{Frames} & \textbf{Duration} & \textbf{Grasp} & \textbf{Release} \\
\midrule
\texttt{lfdws\_t004} & 1745 & 116.8\,s & $t{=}39.47$\,s & $t{=}54.66$\,s \\
\texttt{lfdws\_t005} & 726  & 49.0\,s  & $t{=}10.69$\,s & $t{=}37.50$\,s \\
\bottomrule
\end{tabular}
\caption{Detected grasp/release events on the two new trials. No press
event on either (no wrench data).}
\end{table}

\subsection{Propagation results}

The auto-seed picker placed a clean, small seed on the fastener at the
grasp frame on both trials (visually confirmed), and bidirectional SAM\,2
propagation (\texttt{Code/propagate\_demo\_bidir.py}) gave full-demo
coverage on both: 726/726 frames on \texttt{lfdws\_t005}, 1745/1745 on
\texttt{lfdws\_t004}. Mask area stayed small and stable throughout
(maximum $\sim\!6{,}000$\,px on both, under 1.5\% of the frame), with no
runaway or background-tracking behaviour. Spot-checked at the seed frame,
mid-demo, and a frame where the fastener is genuinely out of view (mask
area correctly reads zero rather than hallucinating a mask). Both trials'
sidecars were composed with \texttt{build\_sidecar.py} in the
grasped-object-only configuration, since neither has a contact-receiver
event to seed from.

\begin{figure}[H]
\centering
\includegraphics[width=0.6\textwidth]{figures/t004/propagation_grasped/f0872_1782993012118220374.png}
\caption{Fastener tracked mid-demo on \texttt{lfdws\_t004}, after being
placed into a mounting hole on the pegboard.}
\label{fig:t004mid}
\end{figure}

\subsection{Updated asks}

The calibration ask stands unchanged. New: confirmation of whether the F/T
sensor was intentionally disconnected for these two recordings, since both
otherwise-complete trials have no contact/press signal at all.

\clearpage

\section{Update --- Calibration Data Arrives, Eye-in-Hand Rig Model}

Dated 10 July. This section covers the lab's response to the asks above:
camera intrinsics, CAD drawings of the camera bracket, and two further
recordings with the force-data issue fixed.

\subsection{The camera is eye-in-hand, not fixed}

The lab's reply clarified the rig geometry: the ZED Mini is mounted on a
bracket bolted to the gripper body, not fixed in the world as earlier
sections assumed. \texttt{current\_pose} is the pose of the \textbf{Bota
SensONE origin} in the base frame --- not the fingertip TCP. Consequences
for the calibration model:
\begin{itemize}[noitemsep]
  \item There is no fixed \texttt{base}$\to$\texttt{camera} transform; it is
        a per-frame quantity, $T_{\text{base}\to\text{camera}}(t) =
        T_{\text{base}\to\text{bota}}(t)\,T_{\text{bota}\to\text{camera}}$,
        with the first term read directly from \texttt{current\_pose} at
        each row.
  \item The wrench is already measured in the bota frame, and
        \texttt{current\_pose} gives that frame's orientation in the base
        frame directly --- so the previously-requested
        \texttt{bota\_frame}$\to$\texttt{base} mount transform is no longer
        needed at all. Rotating the wrench into the base frame uses the same
        rotation as the camera projection.
  \item The one remaining fixed unknown is $T_{\text{bota}\to\text{camera}}$
        --- rigid bracket geometry, obtainable from CAD.
\end{itemize}

\subsection{Camera intrinsics}

Real intrinsics were received (\texttt{intrinsics.md}) for three ZED
resolution modes. The mode matching this pipeline's exported frame size
(the ZED ROS wrapper's 960$\times$540 downscale of HD1080) was written into
\texttt{calibration.yaml} and marked \texttt{filled: true}:
$f_x=f_y=752.9597$, $c_x=472.6555$, $c_y=271.9621$, verified against the
actual PNG dimensions on disk for both new recordings below.

\subsection{CAD extraction attempt for $T_{\text{bota}\to\text{camera}}$}

The lab has not yet run a calibration procedure for the bracket; instead
CAD drawings and a STEP assembly file of the camera bracket were provided,
with the lab separately arranging for someone to extract the transform
from them. In parallel, three independent extraction attempts were made:

\begin{enumerate}[noitemsep]
  \item \textbf{Assembly-tree placement math.} A STEP-file parser
        (\texttt{Code/cad\_extract\_transform.py}) walks the
        \texttt{NEXT\_ASSEMBLY\_USAGE\_OCCURRENCE} + \texttt{AXIS2\_PLACE}-
        \texttt{MENT\_3D} chain to compose each component's rigid placement
        to the assembly root. This recovered the camera housing's origin and
        orientation exactly: translation $[-8.70,-198.46,-25.89]$\,mm in the
        assembly root frame, a $45^\circ$ rotation about the bracket-arm
        axis. Cross-checked two ways: against the PDF drawing's independent
        dimensions ($151.00/177.16/158.82/91.72$\,mm, $30.01^\circ$ tilt),
        and against FreeCAD's own Placement readout for the same part
        (opened interactively) --- both matched exactly.
  \item \textbf{Direct lens-bore search.} Using \texttt{pythonocc-core}
        (OpenCascade's Python bindings) via
        \texttt{Code/cad\_find\_lens\_occ.py}, every cylindrical face in the
        assembly was enumerated and filtered to a plausible lens-radius
        range. Two candidate bores were found near the housing, but they
        are $88$\,mm apart, not the ZED Mini's published $63$\,mm stereo
        baseline --- almost certainly mounting/dowel holes for the housing
        shell, not the lens barrels. The lens glass may be modelled as a
        flat disc rather than a cylindrical bore, or not separately
        labelled; this search was abandoned rather than force an
        unsupported number.
  \item \textbf{CAD orientation + published baseline.} The housing's local
        $X$-axis (from the assembly-tree extraction) is exactly aligned with
        the assembly root's global $X$-axis --- the camera's $45^\circ$ tilt
        is entirely in its $Y/Z$, so the stereo baseline stays level, as
        expected for a camera that pitches down at a workspace while
        keeping its two lenses side by side. Combining this confirmed
        direction with the ZED Mini's published $63$\,mm baseline gives two
        candidate lens positions (bota frame, mm): $A=[111.3,-143.6,25.9]$,
        $B=[155.9,-99.0,25.9]$.
\end{enumerate}

This value is written into \texttt{calibration.yaml} as a documented
\textbf{preliminary} estimate (still \texttt{filled: false}): the housing
pose is solid (three independent cross-checks agree), but which of $A/B$ is
the true left lens is not visually confirmed, and a $\sim$1--2\,cm
lens-to-body depth offset is assumed zero. Good enough to develop and test
against; not yet trusted for sub-centimetre wrench-ray accuracy.

\subsection{\texttt{project\_ee.py} reworked for eye-in-hand}

The projection script was rewritten around the per-frame extrinsic model
above: the anchor point (bota origin, standing in for the end-effector until
a fingertip-TCP offset is available) and the Bicchi wrench ray are computed
in the bota frame and carried into the base frame via the same
\texttt{current\_pose} rotation used for the camera, then projected through
$T_{\text{base}\to\text{camera}}(t)$ built fresh at each event.

Verified in DRY mode (calibration not yet trusted, so nothing is drawn)
across all four sensor-availability profiles seen so far --- full sensors,
force-only, gripper-only, and the original trial as a regression check ---
confirming identical wrench/force numbers to the previously-documented
values on \texttt{lfdws\_t001} ($|F|=26.37$\,N at the press, matching
\S3's Table~2 exactly). The LIVE code path was also exercised with the
calibration flag temporarily set, confirming the transform composition and
image projection execute correctly end-to-end; as expected given the
preliminary lens position, the projected wrench ray did not land on the
contact point, and the flag was reverted to \texttt{filled: false}
afterward. One incidental sanity check: with the camera eye-in-hand, the
anchor point projects to the identical pixel across every event regardless
of where the arm is in the world, since the bota origin sits at a fixed
offset from the camera at all times --- exactly the expected eye-in-hand
invariant, and a useful confirmation that the per-frame transform is wired
correctly.

\subsection{Two more recordings, force data fixed}

Two further recordings arrived, extracted the same way, both with the
gripper, wrench, and depth topics all live simultaneously for the first
time. One is a short test clip (9.8\,s); the other (65.5\,s) is a full task:
a Rubik's cube grasped and pressed down onto the table repeatedly. Unlike
the earlier plate/screwdriver trial, the strongest contact event here is
the cube touching the plain table surface, not a second distinct object ---
recognised as such by inspecting the event frame directly rather than
seeding a meaningless table-texture mask, and the trial was processed in
the grasped-object-only configuration as a result.

The first propagation attempt, seeded with a single point prompt, only
recovered one sticker face of the cube rather than the whole object ---
SAM's point-prompt convention segments the locally-contiguous coloured
region, not the physical object boundary, on a multi-coloured object like a
Rubik's cube. Reseeding with a box prompt around the full cube instead
(\texttt{Code/propagate\_object\_n.py} extended to accept
\texttt{--seed\_box} as an alternative to a point) fixed this: the
corrected propagation tracks the whole cube across all 975 frames, mask
area topping out around 15\% of the frame with no drift.

\end{document}
